\section{Conclusion}
This revision re-examined MultiCardioFusion, an explainable pipeline that enriches the 13 Cleveland clinical variables with 23 ECG-proxy features and benchmarks a parallel PTB-XL ECG branch, under a controlled ablation protocol of \nVariants{} configurations. The system discriminates CAD well (AUC \fullAUC{} \fullAUCci{} under leak-free repeated CV; up to \extHunFu{} at an external hospital) and its explanations are stable for the dominant predictors. The central v4 claim, however, does not survive: the ECG-proxy block does not improve on the clinical variables alone ($\Delta$AUC \dClinAUC{} in favour of clinical-only). It behaves like added noise in control experiments, does not help under 16 perturbations, in any subgroup or at any external site, and the reported $+7.04$~pp gain arose from comparing a favourable hold-out split (95th percentile of \hoN{} splits) with a cross-validation estimate. SMOTE, the specific imbalance strategy and cross-attention over a length-1 sequence likewise add nothing measurable, while the availability of \texttt{ca} and \texttt{thal}, and recalibration of the predicted probabilities, matter.

The contribution of this work is therefore a reproducible, leakage-aware and fully ablated baseline for explainable CAD prediction on Cleveland, together with a demonstration of how single-split evaluation, pre-CV oversampling and attribution to engineered features can produce misleading conclusions on small clinical datasets. Testing true clinical--ECG fusion requires a cohort in which both modalities come from the same patients; the protocol and code released with this paper are designed to be applied to such data directly.
